\section*{Bab \ref{ch:b3:statistical-thermo-applied} --- Penerapan Termodinamika Statistik}

\begin{solution}{exo:b3:statistical-thermo-applied:1}
$\frac32R + R + (3 \times 3 - 5)R = 6.5R = \qty{54.0}{J.K^{-1}.mol^{-1}}$. Pada suhu
kamar, translasi dan rotasi bersifat klasik, kedua vibrasi ulur membeku, dan
vibrasi tekuk yang terdegenerasi rangkap dua, vibrasi yang paling rendah,
tereksitasi sebagian.
\end{solution}

\begin{solution}{exo:b3:statistical-thermo-applied:2}
$x = 1$: $\eu/(\eu - 1)^2 = 0.921$. $x = 10$: $100\eu^{10}/(\eu^{10} - 1)^2 \approx 100\eu^{-10} = \num{4.5e-3}$.
\end{solution}

\begin{solution}{exo:b3:statistical-thermo-applied:3}
\qty{20}{K}: $x_{\mathrm{para}} = 1/(1 + 3 \times 3\eu^{-2 \times 85.35/20}) = 1/(1 + 9 \times \num{1.96e-4}) = 0.998$.
\qty{77}{K}: jumlah genap $1 + 5\eu^{-6.651} = 1.0065$; jumlah ganjil
$3(3\eu^{-2.217} + 7\eu^{-13.30}) = 0.9806$;
$x_{\mathrm{para}} = 1.0065/1.9871 = 0.51$.
\end{solution}

\begin{solution}{exo:b3:statistical-thermo-applied:4}
Untuk $I = 1$ terdapat $3 \times 2 = 6$ keadaan spin simetris dan $3 \times 1 = 3$
keadaan antisimetris; deuteron adalah boson, sehingga fungsi gelombang
totalnya simetris: keadaan spin simetris dengan $J$ genap (orto, bobot 6),
keadaan antisimetris dengan $J$ ganjil (para, bobot 3). Ortodeuterium, yang
memuat $J = 0$, stabil pada 0 K; pada suhu kamar orto:para = 2:1.
\end{solution}

\begin{solution}{exo:b3:statistical-thermo-applied:5}
$\Delta_rE_0 = 2 \times 1883.7 - 2170.3 - 1542.3 = \qty{54.8}{cm^{-1}}$; $\eu^{-54.8/207.2} = 0.768$,
dan $4 \times 0.768 = 3.07$. Nilai lengkapnya, 3.26, 6\,\% lebih tinggi: faktor
translasi, rotasi, dan vibrasi hanya saling meniadakan dengan tepat dalam
limit klasik, dan rotasi molekul-molekul ringan ini masih sebagian
terkuantisasi pada \qty{298}{K}.
\end{solution}

\begin{solution}{exo:b3:statistical-thermo-applied:6}
\ce{N2O}: dua orientasi, $R\ln2 = \qty{5.76}{J.K^{-1}.mol^{-1}}$. \ce{CH3D}: empat,
$R\ln4 =
\qty{11.53}{J.K^{-1}.mol^{-1}}$ (batas atas, yang tercapai jika susunannya acak).
\end{solution}

\begin{solution}{exo:b3:statistical-thermo-applied:7}
$x = 797.6/300 = 2.659$: fungsi Einstein $x^2\eu^{-x}/(1 - \eu^{-x})^2 = 0.572$;
$C_V/R = 3.072$, $C_p = 4.072R = \qty{33.86}{J.K^{-1}.mol^{-1}}$, 0.4\,\% di bawah
nilai tabel (anharmonisitas).
\end{solution}

\begin{solution}{exo:b3:statistical-thermo-applied:8}
$x = 1.029$: $x^2\eu^{-x}/(1 - \eu^{-x})^2 = 0.916$, yaitu \qty{7.62}{J.K^{-1}.mol^{-1}},
92\,\% dari $R$: vibrasi iodin hampir klasik pada suhu kamar.
\end{solution}

\begin{solution}{exo:b3:statistical-thermo-applied:9}
B lebih kaya \ce{^{18}O}. $\alpha_{\mathrm{A/B}} = (1 - 0.0100)/(1 + 0.0020) = 0.98802$.
\end{solution}

\begin{solution}{exo:b3:statistical-thermo-applied:10}
$\Delta\mathrm{ZPE} = \frac12(1 - 1/\sqrt2)(2000 - 3000) = \qty{-146}{cm^{-1}}$;
$K = \eu^{146.4/207.2} = 2.0$. Deuterium lebih suka pergi ke X, ikatan yang
lebih kaku, tempat penghematan \omterm{def:b3:quantum-model-systems:oscillator}{energi titik nolnya} paling besar.
\end{solution}

\begin{solution}{exo:b3:statistical-thermo-applied:11}
$\Delta_rH^\circ = R\ln10\,(-1.766 + 3.392)/(1/900 - 1/1100) = 19.145 \times 1.626/\num{2.020e-4} =
\qty{154}{kJ/mol}$. Kedua atom membawa entalpi $2 \times \frac52RT$; molekulnya
$\frac52RT + RT$ (rotasi) ditambah energi vibrasinya
$R\theta_{\mathrm V}/(\eu^{\theta_{\mathrm V}/T} - 1)$, sedikit kurang dari $RT$: selisihnya,
sekitar \qty{5}{kJ/mol} pada \qty{1000}{K}, ditambahkan pada $D_0$.
\end{solution}

\begin{solution}{exo:b3:statistical-thermo-applied:12}
Tingkat 0 ($g = 1$) dan $6hcB$ ($g = 5$): $q = 1 + 5\eu^{-x}$,
$\langle\varepsilon\rangle/kT = 5x\eu^{-x}/q$, $\langle\varepsilon^2\rangle/(kT)^2 =
5x^2\eu^{-x}/q$, dan $C/R = \langle\varepsilon^2\rangle/(kT)^2 - (\langle\varepsilon\rangle/kT)^2 = 5x^2\eu^{-x}/(1 + 5\eu^{-x})^2$.
Pada \qty{100}{K}, $x = 5.121$: $C/R = 0.783/1.061 = 0.738$.
\end{solution}

\begin{solution}{pb:b3:statistical-thermo-applied:1}
\textbf{1.} $2 \times 2 = 4$: tiga simetris (triplet), satu antisimetris (singlet).
\textbf{2.} Proton adalah fermion, sehingga fungsi gelombang totalnya berganti
tanda pada pertukaran; fungsi rotasi memberikan $(-1)^J$: $J$ genap menyertai
singlet yang antisimetris (para), dan $J$ ganjil menyertai triplet yang
simetris (orto).
\textbf{3.} Tingkat genap 1, tingkat ganjil 3.
\textbf{4.} Pada \qty{300}{K} ($T = 3.5\,\theta_{\mathrm R}$), jumlah rotasi genap dan
ganjil hampir sama; dengan bobot 1 dan 3 keduanya memberikan 1:3 (fraksi para
0.2507).
\textbf{5.} $F(1) = 2 \times 59.322 - 4 \times 0.0471 = \qty{118.46}{cm^{-1}} = \qty{1.417}{kJ/mol}$.
\textbf{6.} $1/(1 + 9\eu^{-170.4/20.37}) = 1/(1 + 9 \times \num{2.3e-4}) = 0.998$.
\textbf{7.} Dengan hanya $J = 0$ dan 1, setengah para berarti
$9\eu^{-2\theta_{\mathrm R}/T} = 1$, $T = 2\theta_{\mathrm R}/\ln9 =
0.91\,\theta_{\mathrm R} = \qty{78}{K}$: persaingannya adalah antara bobot 9 dari $J = 1$
dan faktor Boltzmann tingkat itu.
\textbf{8.} 0.75. Tanpa katalis, perubahan itu memerlukan berhari-hari,
sedangkan pencairan berjam-jam.
\textbf{9.} $(0.75 - 0.002) \times 1.417 = \qty{1.060}{kJ/mol}$.
\textbf{10.} $1.060/\num{2.016e-3} = \qty{526}{kJ/kg}$.
\textbf{11.} $0.905/\num{2.016e-3} = \qty{449}{kJ/kg}$.
\textbf{12.} Kalor perubahan itu 1.17 kali entalpi penguapannya.
\textbf{13.} Kalor itu dihasilkan di dalam cairan, tidak dibawa masuk melalui
dinding.
\textbf{14.} $x_{\mathrm o}(t) = x_{\mathrm{eq}} + (0.75 - x_{\mathrm{eq}})\eu^{-kt}$, dengan
$x_{\mathrm{eq}} = 0.002$.
\textbf{15.} $0.002 + 0.748\eu^{-0.274} = 0.571$.
\textbf{16.} $(0.75 - 0.571) \times 1.417 = \qty{0.254}{kJ/mol}$.
\textbf{17.} $0.254/0.905 = 0.28$: 28\,\% isi tangki dalam sehari.
\textbf{18.} Setelah \qty{168}{h}, $x_{\mathrm o} = 0.112$, kalornya \qty{0.904}{kJ/mol},
sama dengan entalpi penguapan: dalam model ini tangki itu kosong. Perkiraan
itu menaksir kehilangan terlalu tinggi karena molekul-molekul yang menguap
membawa pergi \omterm{def:b3:statistical-thermo-applied:spin-isomers}{hidrogen orto} yang belum berubah; kehilangan yang sebenarnya
lebih kecil tetapi tetap besar.
\textbf{19.} $t_{1/2} = \ln2/0.0114 = \qty{61}{h}$.
\textbf{20.} Di dalam pencair, pada unggun-unggun katalis pada suhu-suhu yang
berurutan, agar mesin pendingin membuang kalor perubahan sebelum cairannya
disimpan.
\textbf{21.} Ya: fraksi para kesetimbangan pada \qty{300}{K} adalah 0.2507.
\textbf{22.} Dalam hidrogen normal, kedua bentuk tidak dapat saling berubah,
sehingga rotasi masing-masing membeku pada tangga tingkatnya sendiri; dalam
hidrogen kesetimbangan, sebagian kalor yang diberikan mengubah para menjadi
orto.
\textbf{23.} Sekitar 3.57 di dekat \qty{50}{K}: selain mengeksitasi rotasi,
kalor itu mengubah molekul dari $J = 0$ ke $J = 1$, \qty{118}{cm^{-1}} lebih
tinggi, sebuah reaksi dengan entalpi positif.
\textbf{24.} 1.50: rotasi membeku pada kedua bentuk (para pada $J = 0$, orto
pada $J = 1$).
\textbf{25.} \textbf{Kalor perubahan / entalpi penguapan $\approx 1.2$ untuk
hidrogen normal}: perubahan itu saja dapat menguapkan seluruh isi tangki,
itulah sebabnya hidrogen diubah menjadi para sebelum disimpan.
\end{solution}
